% Generated draft: requires mathematical and source review.
\documentclass[11pt]{article}
\usepackage[margin=25mm]{geometry}
\usepackage{fontspec}
\setmainfont{Latin Modern Roman}
\usepackage{amsmath,amssymb,mathtools,unicode-math}
\setmathfont{Latin Modern Math}
\usepackage{enumitem,needspace,xurl,hyperref}
\setlength{\parindent}{0pt}
\setlength{\parskip}{5pt}
\setlength{\emergencystretch}{3em}
\newcommand{\N}{\mathbb{N}}
\newcommand{\Z}{\mathbb{Z}}
\newcommand{\Q}{\mathbb{Q}}
\newcommand{\R}{\mathbb{R}}
\newcommand{\C}{\mathbb{C}}
\newcommand{\F}{\mathbb{F}}
\newcommand{\A}{\mathbb{A}}
\newcommand{\PP}{\mathbb{P}}
\newcommand{\E}{\mathbb{E}}
\newcommand{\eps}{\varepsilon}
\newcommand{\dd}{\,\mathrm{d}}
\begin{document}
\section*{UnsolvedMath ID: 20003115}
\textbf{Source problem code:} catalogue entry
\subsection{Definitions and mathematical statement}
M. Kapovich asks for an infinitely generated Fuchsian group $\Gamma_0$ of the first kind whose critical exponent is nonconstant on its quasiconformal Teichm\"uller space $\mathcal T_{qc}(\Gamma_0)$. First kind means its limit set is the full ideal circle. For a discrete group acting on $\mathbb H^2$, the critical exponent is
\[
\delta(\Gamma)=\inf\left\{s>0:\sum_{\gamma\in\Gamma}e^{-s d(o,\gamma o)}<\infty\right\},
\]
independent of $o$. Full limit set and critical exponent $1$ are not equivalent for arbitrary infinitely generated groups.

The aid's all-Fuchsian construction starts with a closed oriented surface $\Sigma_g$, $g\geq2$, and
\[
\pi_1(\Sigma_g)=\langle a_1,b_1,\ldots,a_g,b_g\mid\prod_i[a_i,b_i]=1\rangle.
\]
Let $F_g=\langle x_1,\ldots,x_g\rangle$ and let $K$ be the kernel of the epimorphism $a_i\mapsto x_i$, $b_i\mapsto1$. For each closed hyperbolic metric $m$, let $\rho_m$ be its marked uniformization and set $\Gamma_m=\rho_m(K)$. The aid asserts that these groups are infinitely generated and of the first kind. The associated regular cover $\mathbb H^2/\Gamma_m\to\Sigma_g$ has deck group $F_g$.

Write $\lambda_0(m)$ for the spectral bottom on the cover, defined by the infimum of $\int|\nabla u|^2/\int|u|^2$ over nonzero compactly supported smooth $u$. The aid invokes Brooks's theorem to obtain $\lambda_0(m)>0$, since $g\geq2$ makes the deck group nonamenable and the base is closed. The Elstrodt--Patterson--Sullivan relation is
\[
\lambda_0(m)=
\begin{cases}
\delta(\Gamma_m)(1-\delta(\Gamma_m)),&\delta(\Gamma_m)\geq\tfrac12,\\
\tfrac14,&\delta(\Gamma_m)\leq\tfrac12.
\end{cases}
\]
In particular $\delta(\Gamma_m)<1$ for every such metric.

For a chosen family pinching the disjoint cut-system curves $b_1,\ldots,b_g$ simultaneously to length $\ell\to0$, with the remaining geometric data controlled, the aid reports a collar-cutoff estimate
\[
0<\lambda_0(m_\ell)\leq\frac{C_g}{|\log\ell|^2}.
\]
Eventually $\lambda_0<1/4$, so the relevant critical exponent is on the branch above $1/2$, and
\[
0<1-\delta(\Gamma_{m_\ell})\leq2\lambda_0(m_\ell)
=O(|\log\ell|^{-2}).
\]
Thus these exponents tend to $1$ while staying strictly below $1$. Lifting quasiconformal changes of the compact base places all these marked covers in $\mathcal T_{qc}(\Gamma_{m_0})$; the aid describes the resulting family as a pullback copy of finite-dimensional Teichm\"uller space. Its exponent is therefore nonconstant, giving the reported affirmative construction. The exact pinching-family and collar-test details require verification; the estimate is not asserted uniformly over every unbounded choice of remaining parameters.

The canonical extraction is cut at a page break and ends with an unfinished sentence about a related question. The supplied corrected statement isolates the nonconstant-exponent construction; missing subsequent context is not invented.
\subsection{Short English statement}
The supplied construction uses a free-deck-group cover of a closed surface. Its spectral bottom is positive, keeping the critical exponent below one, but pinching a cut system drives the spectral bottom to zero and the exponent to one within the same quasiconformal deformation space.
\subsection{Sources}
\begin{itemize}
\item Local UnsolvedMath registry, permanent ID $20003115$, source code \texttt{catalogue entry}; canonical \texttt{aim-topology-notes.json}, \texttt{notes[202]}.
\item AIM \emph{Problems on Thurston Metric}, Problem 4.3, attributed to M. Kapovich, workshop \emph{Lipschitz metric on Teichmueller space}: \url{https://aimath.org/pastworkshops/lipschitzteichproblems.pdf}. The supplied verification reports the page-break truncation and corrected statement; status is \texttt{corrected\_verified}. Original remarks and literature field are empty.
\item The aid invokes Brooks's spectral/nonamenability theorem and the Elstrodt--Patterson--Sullivan spectral formula without authoritative full bibliographic references.
\item Research attempt 1: \texttt{full\_solution}, status at run \texttt{unknown}, current status \texttt{solved}, classification \texttt{SOLVED-BY-YOU}, rigor \texttt{proved}, novelty confidence low. The classification is a supplied label, not current authorship or independent review. The machine-generated aid is not peer reviewed and requires independent verification before citation.
\end{itemize}
\end{document}
